% =====================================================================
%  CHAPTER 4: 儿童少年心理行为发育（对应大纲第五章：儿童心理发育）
%  参考往届笔记：第四章 儿童少年心理行为发育
% =====================================================================
\chapter{儿童少年心理行为发育}
\setcounter{chapter}{4}
\chapterintro{
\textbf{大纲要求：}
\begin{itemize}[leftmargin=*, itemsep=2pt]
    \item \textbf{掌握}（双横线）：儿童心理发育的特点；儿童少年认知能力发育
    \item \textbf{熟悉}（单横线）：脑发育（脑发育的关键期、阶段性特征、影响因素及与心理行为发育的关系）；儿童少年情绪情感发育；儿童少年个性及社会化发展
\end{itemize}
\vspace{4pt}
本章介绍儿童少年心理行为发育的全过程。心理行为发育包括认知、语言、情绪情感、人格和社会化适应等多个方面，其发展趋势是从简单到复杂、从低级到高级的上升过程。
}

% ----- 5.4 脑发育 -----
\section{脑发育}

\begin{keybox}
\textbf{\imp{脑发育}}

\textbf{一、脑发育的关键期}

生命早期（受精卵形成-出生后2年），大脑以惊人的速度生长。胎儿期的最后3个月（孕28周）和婴儿出生后2年被称作"大脑发育加速期"，因为成人脑一半以上（75\%）的重量都在该阶段获得。这一时期是脑发育最关键的"机会窗口"，脑的结构和功能具有高度的可塑性，适宜的环境刺激和充足的营养供给对脑的正常发育至关重要。

\textbf{二、脑发育的阶段性特征}

脑的发育遵循严格的程序性规律，主要包括以下阶段：

\begin{enumerate}[leftmargin=*]
    \item \imp{神经元增殖（Proliferation）}：胚胎期神经管形成后，神经上皮细胞迅速分裂增殖，产生大量的神经母细胞。神经元增殖的高峰期在孕10-26周，此阶段每天可产生数十万个神经元。

    \item \imp{神经元迁移（Migration）}：新生成的神经元从室管膜区沿放射状胶质细胞纤维向外迁移，到达大脑皮质的特定位置。迁移过程从内向外，最早迁移的神经元位于皮质深层，后迁移的神经元穿过已定位的细胞层到达更浅表的位置。

    \item \imp{神经元分化（Differentiation）}：神经元到达最终位置后，开始分化形成特定的形态和功能特征，包括轴突和树突的生长、特定神经递质表型的确定等。

    \item \imp{突触形成（Synaptogenesis）}：神经元之间建立突触连接。突触形成在出生前后加速，婴幼儿期达到高峰。2-3岁时大脑的突触密度达到成人水平的150\%，随后经历"突触修剪（Synaptic Pruning）"过程，使神经回路更加高效和精确。

    \item \imp{髓鞘化（Myelination）}：少突胶质细胞包裹轴突形成髓鞘，加速神经冲动的传导。髓鞘化从胎儿期开始，持续至青春期甚至成年早期，遵循从低级中枢到高级中枢、从感觉区到运动区再到联合区的顺序。
\end{enumerate}

\textbf{三、脑发育的影响因素}

\begin{enumerate}[leftmargin=*]
    \item \imp{营养因素}：蛋白质、必需脂肪酸（尤其是DHA）、铁、锌、碘等微量营养素的充足供给是脑正常发育的物质基础。孕期和婴幼儿期营养不良可导致脑细胞数量减少、突触连接异常、髓鞘化延迟，造成不可逆的认知功能损害。

    \item \imp{环境刺激}：丰富的环境刺激（如语言交流、游戏、探索活动等）可促进突触形成和神经回路的精细化。早期感觉刺激的剥夺（如视觉剥夺）可导致相应皮质区域的功能发育障碍。

    \item \imp{应激与压力}：慢性应激可导致糖皮质激素水平持续升高，损害海马神经元，影响学习和记忆功能。安全、稳定的依恋关系对婴幼儿脑发育具有保护作用。

    \item \imp{遗传与表观遗传}：遗传因素决定了脑发育的基本框架，但环境因素可通过表观遗传机制（DNA甲基化、组蛋白修饰等）影响基因表达，调节脑的发育轨迹。

    \item \imp{疾病与毒素}：孕期感染（如风疹病毒、弓形虫）、酒精、烟草、铅、汞等神经毒素可严重干扰脑的正常发育。
\end{enumerate}

\textbf{四、脑发育与心理行为发育的关系}

脑是心理行为活动的物质基础。脑的发育为心理行为发育提供了神经生物学前提：

\begin{itemize}[leftmargin=*]
    \item 前额叶皮质（执行功能、注意控制、计划决策）的缓慢成熟（持续至25岁左右）决定了儿童少年自我控制能力和抽象思维能力的渐进发展。
    \item 边缘系统（情绪加工）的早期成熟与前额叶皮质控制功能发展滞后之间的不平衡，是青春期情绪波动和冒险行为的神经基础。
    \item 语言区（Broca区和Wernicke区）的髓鞘化和突触修剪与儿童语言关键期相对应。
    \item 脑发育的阶段性规律解释了儿童认知能力（如感知觉、注意、记忆、思维）的年龄特征。
    \item 早期脑发育异常（如孤独症谱系障碍、注意缺陷多动障碍等）可导致相应的心理行为问题。
\end{itemize}
\end{keybox}

% ----- 6.0 心理发育概述 -----
\section{儿童心理发育的特点}

\begin{keybox}
\textbf{心理行为发育（Psychological and Behavior Development）：}从受精卵开始到出生、成熟，儿童心理和行为特征的变化过程，包括认知、语言、情绪情感、人格和社会化适应等多个方面。

\textbf{儿童心理发育的特点：}在整个发育过程中，心理行为发育的总趋势是从简单到复杂、从低级到高级的上升过程，各种心理过程和行为特征相互联系、相互影响、共同发展。
\end{keybox}

% ----- 6.1 认知能力发育 -----
\section{儿童少年认知能力发育}

\begin{keybox}
\textbf{认知（Cognition）：}认识活动或认识过程，是大脑反映客观事物的特征、状态及其相互关系、并揭示事物对人的意义与作用的一种高级心理活动。
\end{keybox}

\subsection{感知觉发育}

\begin{keybox}
\textbf{感知觉（Sensory Perception）：}人脑对当前作用于感觉器官的客观事物的反映。

\textbf{一、感觉（Sense）：}对事物个别属性的反映，依赖于个别感觉器官的活动，分为视觉、听觉、嗅觉、味觉和触觉。

\begin{enumerate}[leftmargin=*]
    \item \textbf{视觉}
    \begin{itemize}[leftmargin=*]
        \item 2个月：视觉集中明显形成（尤其是人脸）。
        \item 4个月：出现辨色视觉。
        \item 6个月以前：视力发展敏感期，视力和立体视觉都逐渐发育。
        \item 2$\sim$3岁：能正确辨别红、黄、绿、蓝4种基本颜色并出现双眼视觉。
        \item 出生后3年内：视觉系统功能基本发育成熟。
    \end{itemize}

    \item \textbf{听觉}
    \begin{itemize}[leftmargin=*]
        \item 出生时：听觉器官发育基本成熟，新生儿能听见声音及区分声音的频率、强度和持续时间。
        \item 2岁以后：听力已经很灵敏，几乎达到成人水平。
        \item 学龄初期儿童：听觉感受性随着年龄增长不断发展，在音调辨别、言语听觉、语音听觉和听觉敏感等方面明显提高。
        \item 青少年时期：听觉能力已基本达到稳定水平。
    \end{itemize}

    \item \textbf{嗅觉}
    \begin{itemize}[leftmargin=*]
        \item 新生儿：嗅觉与成人一样敏感，出生时嗅觉发育已比较成熟，能表现出对不同气味的反应。
        \item 7$\sim$8个月：嗅觉逐渐灵敏，能分辨出芳香气味。
        \item 2岁左右：已能很好辨别各种气味。
    \end{itemize}

    \item \textbf{味觉}
    \begin{itemize}[leftmargin=*]
        \item 味觉在儿童时期最发达，以后逐渐衰退。
        \item 2h的新生儿：能分辨甜、酸、苦、咸等多种味道。
        \item 6个月$\sim$1岁：幼儿在这一阶段味觉发展最灵敏。
    \end{itemize}

    \item \textbf{触觉}
    \begin{itemize}[leftmargin=*]
        \item 触觉是人体发展最早、最基本的感觉；触觉在胎儿期已开始发育，到新生儿阶段已高度敏感，触觉是婴儿认识世界的主要手段之一；对婴儿的抚触就是通过对触觉的刺激，增强婴儿触觉的敏感性，加强对外界的反应，促进其发育。
        \item 2$\sim$3岁：能很好辨别各种物体的不同属性。
        \item 5$\sim$6岁：皮肤觉分辨的精细度逐渐提高，在学龄前期已逐步趋于完善。
    \end{itemize}
\end{enumerate}

\textbf{二、知觉（Perception）：}人对事物整体属性的综合反映，是在感觉基础上发展起来的，依赖多种感觉器官，分为时间知觉、空间知觉和运动知觉。

\begin{itemize}[leftmargin=*]
    \item 婴儿3个月时有了分辨简单形状的能力，6个月以前能辨别大小，9$\sim$12个月能知觉形状。
    \item 随着感觉功能的完善促使幼儿的感知觉进一步发展，主要表现为在空间知觉上辨别形状的能力逐渐增强。
    \item 学龄期儿童的视觉感受性和听觉感受性会随年龄增长而不断发展，空间知觉、方位知觉、距离知觉和时间知觉都有很大进步。
\end{itemize}
\end{keybox}

\subsection{注意力和记忆力}

\begin{keybox}
\textbf{一、注意（Attention）：}心理活动对一定对象有选择的集中，是认知过程的开始，分为无意注意和有意注意。

\begin{itemize}[leftmargin=*]
    \item 学龄期儿童：有意注意进一步发展，注意的稳定性随年龄增长而增强，注意集中时间也随之提高。
    \item 青少年期：注意的集中性和稳定性不断增长，稳定性大概能维持40min。
\end{itemize}

\textbf{二、记忆（Memory）：}人脑对过去经验的反映，包括识记、保持和再现（回忆和再认），需经历感知觉记忆、短时记忆和长时记忆3个阶段。

\begin{itemize}[leftmargin=*]
    \item 婴幼儿时期：记忆迅速发展的第一个时期，条件反射的出现是记忆发生的标志。记忆容量增加显著，但记忆以无意识记为主，有意识记成分在逐渐增加。
    \item 学龄期：由于系统学习需要记住大量东西，此时记忆发生质的变化。
    \item 青少年期：记忆的整体水平处于人生的最佳时期。有意识记日益占主导地位，对抽象材料的记忆能力也明显增强。
\end{itemize}
\end{keybox}

\subsection{思维和想象能力}

\begin{keybox}
\textbf{一、思维（Thinking）：}人脑对客观事物间接、概括的反映，能认识事物的本质和事物间的内在联系，属认知的高级阶段。

\begin{itemize}[leftmargin=*]
    \item 婴幼儿期：思维依靠感知觉和动作完成；思维主要是直觉行动思维，是指婴幼儿在动作中进行的思维。婴幼儿在进行这种思维时，只能反映自身动作所触及的具体事物，而不能离开动作，在感知和动作外思考。
    \item 学龄前期：直观行动思维$\to$具体形象思维$\to$抽象逻辑思维，其中占主导地位的是具体形象思维。
    \item 学龄期：思维的基本特点是从具体形象思维逐步过渡到抽象逻辑思维，这是思维发展过程中质的变化。
    \item 青春期：思维能力有很大发展，抽象逻辑思维处于优势地位，能运用抽象思维突破心理运算的界限和理解各种抽象概念，并获得更多增长新知识的机会。
\end{itemize}

\textbf{二、想象（Imagine）：}对已有表象进行加工改造，形成新形象的过程。思维是想象的基础。

\begin{itemize}[leftmargin=*]
    \item 婴幼儿期：1$\sim$2岁的婴幼儿出现想象的萌芽，主要通过动作和口头言语表达出来；3岁时，随着经验与言语的发育，逐渐产生了具有简单想象的游戏。
    \item 学龄前期：已具有丰富的想象力，集中表现在游戏活动中，非游戏时想象水平较低。
    \item 学龄期：想象的有意性、目的性逐渐增强，创造性成分逐渐增大，现实性逐渐提高。
    \item 青少年期：想象能力随教学活动的深入进一步发展，主要表现为有意想象占主要地位，创造想象日益占优势地位，再创想象能力更加独立、概括和精确，想象内容更加复杂，抽象概括性和逻辑性更高。
\end{itemize}
\end{keybox}

\subsection{语言能力发展}

\begin{keybox}
\textbf{语言（Language）：}人类社会中客观存在的现象，是一种社会上约定俗成的符号系统，是由词汇（包括音形义）按一定语法构成的。

\begin{enumerate}[leftmargin=*]
    \item \textbf{前语言阶段（Prespeech Stage）：}从婴儿出生到第一个有真正意义的词产生前的这一时期。
    \item \textbf{语音的发展：}学龄前儿童发音的正确率随年龄的增长而提高，对韵母的发音比较容易掌握，发音错误多出现于辅音中的翘舌音和齿音。
\end{enumerate}
\end{keybox}

% ----- 6.2 情绪情感发育 -----
\section{儿童少年情绪情感发育}

\begin{defbox}
\textbf{一、概念}

\textbf{情绪（Emotion）：}客观事物是否符合人的需要产生的态度体验，反映客观事物与人的需要间的关系。

\textbf{情感（Emotional Feeling）：}与人的高级社会性需要满足与否相联系的态度体验，是在社会交往的实践中逐渐形成的，如友谊感、道德感、美感和理智感等，是人类独有的一种情绪状态。

\textbf{二、儿童青少年情绪情感发育特点}

\begin{enumerate}[leftmargin=*]
    \item \textbf{婴幼儿情绪情感发育特点：}
    \begin{itemize}[leftmargin=*]
        \item 与生理需求是否得到满足直接相关。
        \item 是儿童与生俱来的遗传本能，有先天性。
    \end{itemize}

    \item \textbf{学龄前儿童情绪情感发育特点：}
    \begin{itemize}[leftmargin=*]
        \item 情绪的冲动性逐渐减少。
        \item 情绪情感以外显性为主，内隐性逐渐增强。
        \item 情绪情感以易变性为主，稳定性逐渐发展。
    \end{itemize}

    \item \textbf{学龄期儿童情绪情感发育特点：}
    \begin{itemize}[leftmargin=*]
        \item 情感内容逐渐丰富。
        \item 情感的深刻性和稳定性逐渐增加。
        \item 高级情感逐渐发展。
    \end{itemize}

    \item \textbf{青少年情绪情感发育特点：}
    \begin{itemize}[leftmargin=*]
        \item 情绪情感迅速而热烈，有两极性。
        \item 内隐性和表现性共同存在。
        \item 高级情感已经有了相当的发展。
    \end{itemize}
\end{enumerate}
\end{defbox}

% ----- 6.3 个性及社会化发展 -----
\section{儿童少年个性及社会化发展}

\subsection{个性的发展}

\begin{defbox}
\textbf{一、个性（Personality）/人格：}个人的整个精神面貌，即具有一定倾向性的心理特征的总和，包括一个人个性倾向性、个性心理特征及自我意识。

\textbf{二、气质（Temperament）：}婴幼儿期出生后最早表现出来的一种较为明显而稳定的个性特征，在婴幼儿社会性发展过程中有重要的地位和作用。
\end{defbox}

\subsection{自我意识的发展}

\begin{defbox}
\textbf{自我意识（Self-Consciousness）：}由自我概念、自我评价和自我（情绪）体验等共同组成的心理过程，也是个体不断社会化的过程。

\begin{itemize}[leftmargin=*]
    \item \textbf{自我概念（Self-Concept）：}个人心目中对自己的印象，包括对自己存在、个人身体能力、性格、态度、思想等方面的认识。
    \item \textbf{自我评价（Self-Evaluation）：}自我意识发展的主要成分和标志。
\end{itemize}
\end{defbox}

\subsection{社会化的发展}

\begin{defbox}
\textbf{社会化（Socialization）：}个体在特定社会文化环境中，形成适合于该社会与文化的人格，掌握该社会公认的行为方式，成为合格社会成员的过程。
\end{defbox}

\newpage
